% =========================================================
% Prenex Normal Form
% =========================================================

\subsection*{Prenex Normal Form}

\begin{remark*}[Section toolkit: Prenex normalization]
Prenex normal form is the predicate-logic analogue of the normal-form
procedures from propositional logic. It is a syntactic normalization procedure,
not a semantic change: each step preserves logical equivalence.
\end{remark*}

\begin{definitionbox}{Definition (Quantifier Prefix)}
\begin{definition}[Quantifier Prefix]
\label{def:quantifier-prefix}
A \emph{quantifier prefix} is a finite string
\[
Q_1x_1\,Q_2x_2\cdots Q_nx_n,
\]
where each \(Q_i\in\{\forall,\exists\}\) and each \(x_i\) is a variable.
\end{definition}
\end{definitionbox}

\begin{remark*}[Standard quantified statement]
\[
\Pi=Q_1x_1\cdots Q_nx_n
\Longleftrightarrow
\forall i\leq n\;(Q_i\in\{\forall,\exists\}\text{ and }x_i\in\mathsf{Var}_{\mathcal L}).
\]
\end{remark*}

\begin{remark*}[Predicate reading]
\(\Pi\in\mathsf{QuantifierPrefix}\) means that \(\Pi\) is a finite string
of quantifier-variable pairs.
\end{remark*}

\begin{remark*}[Interpretation]
A prefix records all quantifier choices before the quantifier-free part of a
prenex formula is read.
\end{remark*}

\begin{dependencies}
\begin{itemize}
  \item \hyperref[def:universal-q]{Universal Formula}
  \item \hyperref[def:existential-q]{Existential Formula}
\end{itemize}
\end{dependencies}

\begin{definitionbox}{Definition (Matrix)}
\begin{definition}[Matrix]
\label{def:matrix-prenex}
In a formula \(Q_1x_1\cdots Q_nx_n\,\psi\), the \emph{matrix} is the trailing
formula \(\psi\). For prenex normal form, the matrix is required to be
quantifier-free.
\end{definition}
\end{definitionbox}

\begin{remark*}[Standard quantified statement]
\[
Q_1x_1\cdots Q_nx_n\,\psi \text{ has matrix } \psi,
\qquad
\psi \text{ quantifier-free in prenex normal form.}
\]
\end{remark*}

\begin{remark*}[Predicate reading]
\(\psi\in\mathsf{Matrix}(\Phi)\) means that \(\psi\) is the matrix of the
formula \(\Phi\).
\end{remark*}

\begin{remark*}[Interpretation]
The matrix is the part of a prenex formula where the propositional connectives
and atomic formulas remain after all quantifiers have been pulled forward.
\end{remark*}

\begin{dependencies}
\begin{itemize}
  \item \hyperref[def:wff-pred]{Well-Formed Formula}
\end{itemize}
\end{dependencies}

\begin{definitionbox}{Definition (Prenex Formula)}
\begin{definition}[Prenex Formula]
\label{def:prenex-formula}
A \emph{prenex formula} is a formula of the form
\[
Q_1x_1\,Q_2x_2\cdots Q_nx_n\,\psi,
\]
where \(Q_1x_1\cdots Q_nx_n\) is a quantifier prefix and \(\psi\) is
quantifier-free.
\end{definition}
\end{definitionbox}

\begin{remark*}[Standard quantified statement]
\[
\Phi \text{ is prenex}
\Longleftrightarrow
\Phi=Q_1x_1\cdots Q_nx_n\,\psi
\text{ with } \psi \text{ quantifier-free.}
\]
\end{remark*}

\begin{remark*}[Predicate reading]
\(\Phi\in\mathsf{PrenexFormula}\) means that all quantifiers of \(\Phi\)
occur in an initial prefix.
\end{remark*}

\begin{remark*}[Negated quantified statement]
\[
\Phi \text{ is not of the form } Q_1x_1\cdots Q_nx_n\,\psi
\text{ with } \psi \text{ quantifier-free.}
\]
\end{remark*}

\begin{remark*}[Interpretation]
Prenex formulas isolate quantifier order from the quantifier-free matrix.
\end{remark*}

\begin{remark*}[Examples]
\(\forall x\,\exists y\,(P(x)\to R(x,y))\) is prenex.
\end{remark*}

\begin{remark*}[Non-Examples]
\(\forall x\,(P(x)\to\exists y\,R(x,y))\) is not prenex because a quantifier
remains inside the matrix.
\end{remark*}

\begin{dependencies}
\begin{itemize}
  \item \hyperref[def:quantifier-prefix]{Quantifier Prefix}
  \item \hyperref[def:matrix-prenex]{Matrix}
\end{itemize}
\end{dependencies}

\begin{definitionbox}{Definition (Prenex Normal Form)}
\begin{definition}[Prenex Normal Form]
\label{def:pnf}
A formula \(\psi\) is a \emph{prenex normal form} of a formula \(\varphi\)
exactly when \(\psi\) is prenex and \(\varphi\equiv\psi\).
\end{definition}
\end{definitionbox}

\begin{remark*}[Standard quantified statement]
\[
\psi \text{ is a prenex normal form of } \varphi
\Longleftrightarrow
\psi \text{ is prenex}
\text{ and }
\varphi\equiv\psi.
\]
\end{remark*}

\begin{remark*}[Predicate reading]
\(\psi\in\mathsf{PrenexNormalForm}(\varphi)\) means that \(\psi\) is a
prenex formula logically equivalent to \(\varphi\).
\end{remark*}

\begin{remark*}[Interpretation]
Prenex normal form changes the placement of quantifiers but not the meaning of
the formula.
\end{remark*}

\begin{remark*}[Exposition]
Prenex normal form is not unique. Different choices of bound-variable names and
different orders of valid pulling steps can produce different but logically
equivalent prenex formulas.
\end{remark*}

\begin{remark*}[Exposition]
A prenex conversion proceeds in stages: eliminate implications and
biconditionals, push negations inward, rename bound variables as needed, pull
quantifiers forward, and leave a quantifier-free matrix.
\end{remark*}

\begin{remark*}[Examples]
Starting with
\[
\neg\bigl(\forall x\,P(x)\to\exists y\,Q(y)\bigr),
\]
eliminate the implication, push the negation inward, and pull the quantifiers
forward to obtain
\[
\forall x\,\forall y\,(P(x)\land\neg Q(y)).
\]
\end{remark*}

\begin{dependencies}
\begin{itemize}
  \item \hyperref[def:prenex-formula]{Prenex Formula}
  \item \hyperref[def:log-equiv]{Logical Equivalence}
\end{itemize}
\end{dependencies}

\begin{theorembox}{Theorem (Existence of Prenex Normal Form)}
\begin{theorem}[Existence of Prenex Normal Form]
\label{thm:pnf}
Every first-order formula is logically equivalent to some formula in prenex
normal form.
\hyperref[prf:pnf]{Go to proof.}
\end{theorem}
\end{theorembox}

\begin{remark*}[Standard quantified statement]
\[
\forall \varphi\in\mathsf{Form}_{\mathcal L}\;
\exists \psi\in\mathsf{Form}_{\mathcal L}\;
(\psi \text{ is prenex}\land \varphi\equiv\psi).
\]
\end{remark*}

\begin{remark*}[Predicate reading]
\(\forall\varphi\,
\exists\psi\,\psi\in\mathsf{PrenexNormalForm}(\varphi)\).
\end{remark*}

\begin{remark*}[Interpretation]
Every formula can be converted into an equivalent one with all quantifiers at
the front and a quantifier-free matrix at the end.
\end{remark*}

\begin{remark*}[Exposition]
The proof proceeds by structural rewriting: remove implication and
biconditional connectives, push negations inward, rename bound variables to
avoid capture, and then pull quantifiers forward using equivalences whose side
conditions have been made true by renaming.
\end{remark*}

\begin{dependencies}
\begin{itemize}
  \item \hyperref[def:pnf]{Prenex Normal Form}
  \item \hyperref[thm:qneg]{Quantifier Negation Laws}
  \item \hyperref[thm:rename]{Renaming Bound Variables}
  \item \hyperref[thm:qdist]{Valid Quantifier Distribution}
\end{itemize}
\end{dependencies}
